\bookStart{Short Lay of Siward}[Sig·urðar kviða in skamma]
\setBookCode{Sigurdskamma}

\begin{flushright}%
\textbf{Dating} \parencite{Sapp2022}: early C11th (0.876)

\textbf{Meter:} \Fornyrdislag%
\end{flushright}

\section{Introduction}

The \textbf{Short Lay of Siward} (\Sigurdskamma) continues the story of Siward, focusing on his death and Byrnhild’s subsequent grieving.  Despite its title, \Sigurdskamma\ has ca. 280 long-lines and is thus one of the longest poems in \Regius, something which suggests that it was once prececed by an even longer poem about Siward (cf.~\textlink{Lacuna} above).

\subsection{Dating}

The language of \Sigurdskamma\ has some archaic and unusual features.

The most notable archaic feature is its use of the particle \emph{of} before nominals.  This particle originally corresponds to an older Germanic \emph{*ga-} (thus Anon \emph{Bjark} 1/4b (\Skp\ 3): \emph{of sinnar} ‘companions’ < PGmc \emph{*ga·sinþaniz} > OE \emph{ge·sïþas} ‘id.’) but becomes obsolete early in ON poetry except for when preceding verbs; in the Scaldic corpus attestations of \emph{of} before nominals occur almost exclusively(?---TODO) before ca. 1000 \ce\ (\Skp\ 5, p. xcic).

In \Sigurdskamma, \emph{of} occurs twice before nouns: 23/1a \emph{of dolgr} (no non-Norse parallels but see note there) and 58/5a: \emph{of skǫp} (= OE \emph{ge·sceap}).  It also occurs twice before the adj.~\emph{líkr}, viz.~in 36/2b and 61/4b: \emph{of líkr} (= OE \emph{ge·líc}).  This usage is unique in ON poetry since the \emph{of} replaces the usual fossilized \emph{g-} in \emph{g·líkr}, which proves that the poet of \Sigurdskamma\ had a strong intuitive sense of the semantics of this particle.  A later archaizer would far more likely have created the hypercorrected form \emph{*of g·líkr}

Another archaic feature is the long form \emph{haufuð} (23/2).  Like the other Walsing-poems (save \textlink{Gripisspa}), \Sigurdskamma\ requires the long first syllable in the name \emph{Sig·vǫrðr} ‘Siward’ \parencite[135--141]{Sievers1889}.

Metrically \Sigurdskamma\ has a notable tolerance for three-syllable verses.

Contrary to Sapp’s statistical model the aforementioned archaic traits allow us to date the poem to the 10th century.  It should therefore be of interest to the reader that the poem’s depictions of widow-suicide (suttee) and the afterlife are likely authentic and pagan.

\sectionline

\newpage

\section{Text}

\subsection{Sig·urðar kviða (‘A Lay of Siward’)}

\bvg\bva%
\edtrans{\edtext{Ár}{\Afootnote{The \emph{A} is an ornamented initial 3 lines in height in \Regius.}} vas þat’s}{It was of yore that}{\Bfootnote{Shared with \textlink{GudrunOne}[1]/1; see there for further parallels.}} \edtext{\alst{S}ig·\emph{vǫ}rðr}{\Afootnote{metr.~emend.~\parencite[135--141]{Sievers1889}; \emph{Sig·urðr} \Regius}} \hld\ \alst{s}ótti Gjúka &
\alst{v}ǫlsungr ungi \hld\ es \alst{v}egit hafði; &
\alst{t}ók við \alst{t}ryggðum \hld\ \alst{t}vęggja brǿðra &
sęldu⸗sk \alst{ęi}ða \hld\ \alst{ę}ljun-frǿknir.\eva

\bvb {\huge I}\textsc{t was of yore} that Siward, the young Walsing \\
who had triumphed, sought out Yivick. \\
He got the truces of two brothers; \\
they exchanged oaths, the men brave of zeal.\evb\evg


\bvg\bva%
\alst{M}ęy buðu hǫ̇num \hld\ ok \alst{m}ęiðma fjǫlð, &
\alst{G}uð·ru̇nu ungu \hld\ \alst{G}júka dóttur; &
\alst{d}rukku ok \alst{d}ǿmðu \hld\ \alst{d}ǿgr mart saman &
\alst{S}ig·\emph{vǫ}rðr ungi \hld\ ok \alst{s}ynir Gjúka,\eva

\bvb They offered him a maiden and a multitude of treasures \\
with young Guthrun, Yivick’s daughter. \\
They drank and discussed for many days and nights together, \\
young Siward and the sons of Yivick,\evb\evg


\bvg\bva%
und’s þęir \alst{B}ryn·hildar \hld\ \alst{b}iðja fóru &
\alst{s}vá’t þęim \alst{S}ig·\emph{vǫ}rðr \hld\ ręið ï \alst{s}inni &
\alst{v}ǫlsungr ungi \hld\ ok \alst{v}ega kunni; &
hann of \alst{ę́}tti \hld\ ef hann \alst{ęi}ga knę́tti.\eva

\bvb till they journeyed to ask for Byrnhild’s hand, \\
whenas Siward rode alongside them, \\
the young Walsing, and he could fight; \\
he would have her if he could have her.\evb\evg


\bvg\bva%
\Ballnote{Siward disguised himself as Guther and married Byrnhild while still in disguise, but out of loyalty for his oath-sworn friend Siward did not sleep with nor even touch her in the bed; instead he placed a sword between them.}%
\alst{S}ęggr inn \alst{s}uðr-ǿni \hld\ lagði \alst{s}verð nøkkvit &
\edtrans{\alst{m}ę́ki \alst{m}ál-fáan}{a picture-painted blade}{\Bfootnote{Also occurring in \textlink{Skirnismal}[23]/1, 25/1; see there.}} \hld\ ȧ \alst{m}eðal þęira &
né hann \alst{k}onu \hld\ \alst{k}yssa gęrði &
né \alst{h}únskr konungr \hld\ \alst{h}ęfja sér af armi &
\alst{m}ęy frum-unga; \hld\ fal hann \alst{m}ęgi Gjúka.\eva

\bvb The southern youth laid a naked sword, \\
a picture-painted blade between them; \\
nor did he kiss the woman, \\
nor did the Hunnish king hold in his arms \\
the most young maiden; he left that to the lad of Yivick.\evb\evg


\bvg\bva%
Hǫ̇n sér at \alst{l}ífi \hld\ \alst{l}ǫst né vissi &
ok at \alst{a}ldr-lagi \hld\ \alst{ę}kki grand &
\alst{v}amm þat’s \alst{v}ę́ri \hld\ eða \alst{v}esa hygði; &
\alst{g}engu þęss ȧ milli \hld\ \alst{g}rimmar urðir.\eva

\bvb She knew no blemish to her life \\
and in her allotted age no wrong, \\
no fault which had been nor which thought to arise; \\
in that the cruel weirds intervened.\evb\evg


\bvg\bva%
\alst{Ęi}n sat hǫ̇n \alst{ú}ti \hld\ \alst{a}ptan dags, &
\edtext{nam hǫ̇n svá bęrt \hld\ umb at mę́la⸗sk:}{\Bfootnote{This line is missing alliteration and no obvious emendation can be found.}} &
\llap{„}Hafa skal’k \alst{S}ig·\emph{vǫ}rð, \hld\ ---eða þó \alst{s}velti!--- &
\alst{m}ǫg frum-ungan, \hld\ \alst{m}ér ȧ armi.\eva

\bvb Alone she sat outside at the close of day; \\
she took to speak plainly thus: \\
“I shall have Siward---else let him die!--- \\
that most young lad in my arms!\evb\evg


\bvg\bva%
\alst{O}rð mę́lta’k nú, \hld\ iðr⸗umk \alst{ę}ptir þęss, &
kvǫ̇n ’s hans \alst{G}uð·ru̇n \hld\ ęn ek \alst{G}unn·ars, &
\alst{l}jótar nornir \hld\ skópu oss \alst{l}anga þrǫ́.“\eva

\bvb The word I just spoke I afterwards regret. \\
Guthrun is his wife, but I am Guther’s; \\
ugly norns shaped for us a long yearning.“\evb\evg


\bvg\bva%
\alst{O}pt gęngr hǫ̇n \alst{i}nnan, \hld\ \alst{i}lls of fylld, &
\edtrans{\alst{í}sa ok \alst{jǫ}kla}{with ice-sheets and icicles}{\Bfootnote{A description of Byrnhild’s inner state; her jealousy makes her so cold that it is as if her insides froze to ice.}}, \hld\ \alst{a}ptan hvęrn, &
es þau \alst{G}uð·ru̇n \hld\ \alst{g}anga ȧ bęð &
ok hana \alst{S}ig·\emph{vǫ}rðr \hld\ \alst{s}vęipr ï ripti, &
\alst{k}onungr inn hu̇nski, \hld\ \alst{k}vǫ̇n frjá sïna.\eva

\bvb She often goes outside filled with hate, \\
ice-sheets, and icicles every evening \\
when he and Guthrun go to their bed \\
and Siward wraps her in the blanket, \\
the Hunnish king his beloved wife.\evb\evg


\bvg\bva%
\llap{„}\alst{V}ǫn gęng’k \alst{v}ilja, \hld\ \alst{v}ers ok \edtext{b\emph{arn}a}{\Afootnote{sens.~emend., cf.~\textlink{GudrunOne}[23]/2; \emph{bęggja} ‘both’ \Regius}}, &
verð’k mik \alst{g}ǿla \hld\ af \alst{g}rimmum hug.“\eva

\bvb “I go without strength of will, husband, or children; \\
I shall have to soothe myself from my cruel thoughts.”\evb\evg


\bvg\bva%
Nam af þęim \alst{h}ęiptum \hld\ \alst{h}vętja⸗sk at vígi: &
\llap{„}Þú skalt, \alst{G}unn·arr, \hld\ \alst{g}ørst of láta &
\alst{m}ïnu landi \hld\ ok \alst{m}ér sjalfri; &
mun’k \alst{u}na \alst{a}ldri \hld\ með \alst{ǫ}ðlingi.\eva

\bvb Out of that hatred she took to incite [men] to war: \\
“Thou shalt, Guther, very soon lose \\
my land and my very self; \\
I will never be content with an athling.\evb\evg


\bvg\bva%
Mun’k \alst{a}ptr fara \hld\ þar’s \alst{á}ðan vas’k &
með \alst{n}á-bornum \hld\ \alst{n}iðjum mïnum, &
þar mun’k \alst{s}itja \hld\ ok \alst{s}ofa lífi &
nema þú \alst{S}ig·\emph{vǫ}rð \hld\ \alst{s}velta látir &
ok \alst{jǫ}furr \alst{ǫ}ðrum \hld\ \alst{ǿ}ðri verðir.\eva

\bvb I will go back whither I was before \\
amidst my close-born kinsmen; \\
there will I sit and sleep my life away \\
unless thou lettest Siward die, \\
and becomest a chieftain nobler than he.\evb\evg


\bvg\bva%
Lǫtum son \alst{f}ara \hld\ \alst{f}ęðr ï sinni; &
skal⸗at \alst{u}lf \alst{a}la \hld\ \alst{u}ngan lęngi. &
\alst{H}vęim verðr \alst{h}ǫlða \hld\ \alst{h}ęfnd lėttari &
---\alst{s}íðr til \alst{s}átta--- \hld\ at \alst{s}onr lifi.“\eva

\bvb We shall let the son follow in his father’s tracks; \\
one shall not rear a young wolf for long. \\
For every man revenge becomes easier \\
---peace less likely---if his son lives.”\evb\evg


\bvg\bva%
\edtext{Ręiðr varð Gunn·arr \hld\ ok hnipnaði}{\Bfootnote{This line is missing alliteration and no obvious emendation can be found.}}, &
\alst{s}vęip \alst{s}ïnum hug, \hld\ \alst{s}at umb allan dag; &
hann \alst{v}issi þat \hld\ \alst{v}il⸗gi gǫr⸗la &
hvat hǫ̇num \alst{v}ę́ri \hld\ \alst{v}inna sǿmst &
eða hǫ̇num \alst{v}ę́ri \hld\ \alst{v}inna bęzt, &
alls sik \alst{V}ǫlsung \hld\ \alst{v}issi firrðan &
ok at \alst{S}ig·\emph{vǫ}rð \hld\ \alst{s}ǫknuð mikinn.\eva

\bvb Then Guther turned wroth and drooped; \\
he twisted his mind, sat for the whole day.  \\
He knew not clearly in any way \\
what was most fitting for him to do, \\
or what was best for him to do, \\
as he knew he would be bereaved of the Walsing \\
and [feel] great loss after Siward.\evb\evg


\bvg\bva%
\alst{Ý}mist hann hugði \hld\ \alst{ja}fn-langa stund, &
þat vas \alst{ęi}gi \hld\ \alst{á}rar títt &
at frȧ \alst{k}onung-dóm \hld\ \alst{k}vȧnir gengi; &
nam hann sér \alst{H}ǫgna \hld\ \alst{h}ęita at ru̇num, &
þar \alst{á}tti hann \hld\ \alst{a}lls full-trúa.\eva

\bvb His thoughts went back and forth for an even-long time. \\
It was not at all desirable in times of yore \\
that wives should depart from a kingdom. \\
He summoned Hawn for private counsel; \\
there he had a friend fully trusted in all.\evb\evg


\bvg\bva%
\speakernote{[Gunn·arr kvað:]}%
\llap{„}\alst{Ęi}n ’s mér Bryn·hildr \hld\ \alst{ǫ}llum bętri, &
of \alst{b}orin \alst{B}uðla, \hld\ hǫ̇n ’s \alst{b}ragr kvenna; &
\alst{f}yrr skal’k mïnu \hld\ \alst{f}jǫrvi láta &
an þęirar \alst{m}eyjar \hld\ \alst{m}ęiðmum tẏna.\eva

\bvb\speakernoteb{[Guther quoth:]}%
“Byrnhild alone is better to me than all--- \\
that woman born of Budle; she is the peak of women. \\
I shall sooner lose my own life \\
than abandon that maiden’s treasures.\evb\evg


\bvg\bva%
Vill þú okkr \alst{f}ylki \hld\ til \alst{f}\emph{éa}r véla? &
Gótt ’s at \alst{r}áða \hld\ \alst{R}ïnar malmi, &
ok \alst{u}nandi \hld\ \alst{au}ði stýra &
ok \alst{s}itjandi \hld\ \alst{s}ę́lu njóta.“\eva

\bvb Wilt thou with me betray the marshaller for money? \\
It is good to rule the ore of the Rhine \ken{gold}, \\
and, prospering, to wield the wealth, \\
and, sitting, to enjoy good fortune.”\evb\evg


\bvg\bva%
\edtrans{\alst{Ęi}nu því Hǫgni \hld\ \alst{a}nn·svǫr vęitti}{This one answer Hawn then granted}{\Bfootnote{Shared with st.~45/1 below and \textlink{Brot}[6]/1.}}: &
\llap{„}\alst{S}amir ęigi okkr \hld\ \alst{s}·líkt at vinna, &
\alst{s}verði rofna \hld\ \alst{s}varna ęiða, &
\alst{ęi}ða svarna, \hld\ \alst{u}nnar tryggðir.\eva

\bvb This one answer Hawn then granted: \\
“It does not befit us two to work such things, \\
to break by the sword the sworn oaths--- \\
the sworn oaths, the taken truces.\evb\evg


\bvg\bva%
Vitum⸗a vit ȧ \alst{m}oldu \hld\ \alst{m}ęnn in sę́lli &
meðan \edtrans{\alst{f}jórir vér}{we four}{\Bfootnote{Guther, Hawn, Godthorm, and Siward.}} \hld\ \alst{f}olki rǫ́ðum &
ok sá inn \alst{h}u̇nski \hld\ \alst{h}ęr-Baldr lifir, &
né in \alst{m}ę́tri \hld\ \alst{m}ę́gð ȧ \alst{m}oldu &
ef vér \alst{f}imm sonu \hld\ \alst{f}ǿðum lęngi &
\alst{ǫ́}ttum góða \hld\ \alst{ǿ}xla knę́ttim.\eva

\bvb We shall not know upon the earth a more blessed man \\
while we four do rule the folk \\
and the Hunnish host-Balder \ken*{\textsc{warrior} = Siward} lives; \\
nor a nobler in-law upon the earth \\
if we rear five sons for long \\
good of lineages, [and] should let them grow.\evb\evg


\bvg\bva%
Ek \alst{v}ęit gǫr⸗la \hld\ hvaðan \alst{v}egir standa: &
Eru \alst{B}ryn·hildar \hld\ \alst{b}rek of·mikil.\eva

\bvb I know clearly whence these troubles arise: \\
Byrnhild’s demands are far too great.\evb\evg


\bvg\bva%
Vit skulum \alst{G}uð·þorm \hld\ \alst{g}ørva at vígi, &
\alst{y}ngra bróður, \hld\ \alst{ȯ}·fróðara; &
hann vas fyr \alst{ú}tan \hld\ \alst{ęi}ða svarna, &
\alst{ęi}ða svarna, \hld\ \alst{u}nnar tryggðir.“\eva

\bvb We shall make Godthorm ready for war, \\
our younger brother, the less learned. \\
He was outside of the sworn oaths, \\
the sworn oaths, the taken truces.”\evb\evg


\bvg\bva%
Dę́lt vas at \edtrans{\alst{ę}ggja \hld\ \alst{ȯ}·bil-gjarnan}{incite the unyielding man}{\Bfootnote{Also found in a proverb in \Grettissaga.}}, &
\edtrans{stóð til \alst{h}jarta \hld\ \alst{h}jǫrr}{the sword was stood unto the heart}{\Bfootnote{Formulaic.  Cf.~\textlink{Voluspa}[52]/3--4, \textlink{Fafnismal}[1]/4.}} Sig·\emph{vę}rði.\eva

\bvb It was easy to incite the unyielding man \ken*{= Godthorm}; \\
the sword was stood unto the heart of Siward.\evb\evg


\bvg\bva%
\Ballnote{Sts.~22--23: Siward throws his sword, Gram, at Godthorm with such strength that it cuts him clean in two parts.}%
Réð til \alst{h}ęfnda \hld\ \alst{h}ęr-gjarn ï sal &
ok \alst{ę}ptir varp \hld\ \alst{ȯ}·bil-gjǫrnum; &
fló til \alst{G}uð·þorms \hld\ \alst{G}rams ramm⸗liga &
\alst{k}yn-birt \emph{éa}rn \hld\ ór \alst{k}onungs hęndi.\eva

\bvb The war-eager one turned to revenge in the hall \\
and threw after the unyielding man; \\
towards Godthorm flew violently Gram’s \\
flashing-white blade from the hand of the king.\evb\evg


\bvg\bva%
\alst{H}né hans \edtrans{of dolgr}{enemy}{\Bfootnote{Identical in sense to \emph{dolgr}.  The prefixing of a noun by the particle \emph{of} is an archaic feature; cf.~Introduction to the poem above.  The authenticity of \emph{of dolgr} is confirmed by another attestation in Egill \emph{Arkv} 22 (\Skp\ 5) from ca. 962 \ce.}} \hld\ til \alst{h}luta tvęggja; &
\alst{h}ęndr ok \edtext{\alst{h}\emph{au}fuð}{\Afootnote{metr. emend.; \emph{hǫfuð} \Regius}\Bfootnote{The long form \emph{haufuð} is metrically required since the second lift cannot fall on a short syllable \parencite[15]{Males2025Of}.}} \hld\ \alst{h}né ȧ annan veg &
ęn \alst{f}óta hlutr \hld\ \alst{f}ell aptr ï stað.\eva

\bvb His enemy sunk down in two parts: \\
his arms and head sunk to the side, \\
but the part of his legs fell backwards to the ground.\evb\evg


\bvg\bva%
\alst{S}ofnuð vas Guð·ru̇n \hld\ ï \alst{s}ę́ingu &
\alst{s}orga-laus \hld\ hjá \alst{S}ig·\emph{vę}rði; &
en hǫ̇n \alst{v}aknaði \hld\ \alst{v}ilja firrð &
es hǫ̇n \edtrans{\alst{F}ręys vinar}{Free’s friend}{\Bfootnote{It is not clear why Siward is called the friend of this god in particular.  TODO: any other kennings of this type?  Cf.~\emph{Ing-wine}.}} \hld\ \alst{f}laut ï dręyra.\eva

\bvb Guthrun lay sleeping in the bed, \\
sorrowless by Siward’s side, \\
but she awoke deprived of strength of will \\
when in the blood of Free’s friend \ken*{= Siward} she swam.\evb\evg


\bvg\bva%
\alst{S}vá \alst{s}ló hǫ̇n \alst{s}várar \hld\ \alst{s}ïnar hęndr &
at \alst{r}amm-hugaðr \hld\ \alst{r}ęis upp við bęð: &
\llap{„}\alst{G}rát⸗a þú, \alst{G}uð·ru̇n, \hld\ svá \alst{g}rimm⸗liga, &
\alst{b}rúðr frum-unga--- \hld\ þér \alst{b}rǿðr lifa.\eva

\bvb So harshly did she beat her hands \\
that the strong-minded man \ken*{= Siward} rose up against the bed: \\
“Weep thou not, Guthrun, so horribly, \\
O most young bride---thy brothers still live!\evb\evg


\bvg\bva%
\Ballnote{Sts.~26--27: Siward predicts the death of his young son, Syemund.}%
\alst{Á}’k til \alst{u}ngan \hld\ \alst{ę}rfi-nytja, &
kann⸗at hann \alst{f}irra⸗sk \hld\ ór \alst{f}jánd-garði; &
þęir \alst{s}ér hafa \hld\ \alst{s}várt ok dátt &%TODO: dátt
ęn \alst{n}ę́r \alst{n}umit \hld\ \alst{n}ý⸗lig rǫ́ð.\eva

\bvb I have an heritance-enjoyer \ken{son} far too young; \\
he cannot remove himself from the court of foes. \\
They make themselves fierce and violent, \\
and have nearly fulfilled their new plans.\evb\evg


\bvg\bva%
Ríðr⸗a þęim \alst{s}íðan \hld\ þó’tt \alst{s}jau alir, &
\alst{s}ystur \alst{s}onr \hld\ \alst{s}·líkr at þingi; &
ek vęit \alst{g}ǫr⸗la \hld\ hví \alst{g}ęgnir nú: &
\alst{ęi}n vęldr Bryn·hildr \hld\ \alst{ǫ}llu bǫlvi.\eva

\bvb With them will not ride---though thou begettest seven--- \\
such a sister’s son to the Thing. \\
I know clearly why this now happens: \\
alone is Byrnhild at fault for all evil.\evb\evg


\bvg\bva%
\Ballnote{Siward comforts Guthrun by blaming Byrnhild for his death while defending the innocence of her brother Guther who has no apparent motive to kill Siward when he has not done anything ill against him.}%
\alst{M}ér unni \alst{m}ę́r \hld\ fyr \alst{m}ann hvęrn &
ęn við \alst{G}unn·ar \hld\ \alst{g}rand ękki vann’k; &
þyrmða’k \alst{s}ifjum, \hld\ \alst{s}vǫrnum ęiðum, &
síðr vę́ra’k \alst{h}ęitinn \hld\ \alst{h}ans kvȧnar vinr.“\eva

\bvb That maiden loved me more than any man, \\
but against Guther I did no wrong; \\
I respected the kinship, the sworn oaths; \\
late will I be called a ‘friend’ of his wife.”\evb\evg


\bvg\bva%
\alst{K}ona varp ǫndu \hld\ ęn \alst{k}onungr fjǫrvi, &
\alst{s}vá sló hǫ̇n \alst{s}vára\emph{n} \hld\ \alst{s}ïnni hęndi &
at \alst{k}vǫ́ðu við \hld\ \alst{k}álkar ï vǫ́ &
\edtrans{ok \alst{g}ullu við \hld\ \alst{g}ę̇ss ï tu̇ni}{and in response shrieked the geese in the yard}{\Bfootnote{Identically shared with \textlink{GudrunOne}[16]/3.}}.\eva

\bvb The woman lost her breath but the king his life. \\
So harshly did she beat her hand \\
that in response clanged the chalices in the cupboard \\
and in response shrieked the geese in the yard.\evb\evg


\bvg\bva%
\edtext{Hló þȧ \alst{B}ryn·hildr, \hld\ \alst{B}uðla dóttir, &
\alst{ęi}nu sinni \hld\ af \alst{ǫ}llum hug}{\lemma{Hló \dots hug ‘Then \dots\ heart’}\Bfootnote{Also found in \textlink{Brot}[9]/1--2 with varation in the first b-verse.}} &
es hǫ̇n til \alst{h}vílu \hld\ \alst{h}ęyra knátti &
\alst{g}jallan \alst{g}rát \hld\ \alst{G}júka dóttur.\eva

\bvb Then laughed Byrnhild, Budle’s daughter, \\
a single time out of all her heart, \\
when from their resting-place she heard \\
the shrill weeping of Yivick’s daughter.\evb\evg


\bvg\bva%
Hitt kvað þȧ \alst{G}unn·arr, \hld\ \alst{g}ramr hauk-stalda: &
\llap{„}\alst{H}lę́r⸗a þú af því, \hld\ \alst{h}ęipt-gjǫrn kona, &
\alst{g}lǫð ȧ \alst{g}olfi, \hld\ at þér \alst{g}óðs viti; &
\alst{h}ví \alst{h}afnar þú \hld\ inum \alst{h}víta lit? &
\alst{F}ęikna \alst{f}ǿðir, \hld\ hygg at \alst{f}ęig sé\emph{i}r.\eva

\bvb This then quoth Guther, the lord of young warriors \\
“Thou dost not laugh, O feud-eager woman, \\
glad on the floor, since thou knowest any good. \\
Why dost thou abandon thy white colour? \\
Rearer of wickedness, I think thou art fey!\evb\evg


\bvg\bva%
Þú \alst{v}ę́rir þęss \hld\ \alst{v}erðust kvinna &
at fyr \alst{au}gum þér \hld\ \alst{A}tla hjǫggim, &
sę́ir \alst{b}rǿðr þïnum \hld\ \alst{b}lóðugt sár, &
\alst{u}ndir dręyrgar, \hld\ knę́ttir \alst{y}fir binda.“\eva

\bvb For this wouldst thou be the worthiest of women \\
that we before thy eyes should cut down Attle, \\
that thou would see on thy brother a bloody gash, \\
gory wounds; then wouldst thou bind over them.”\evb\evg


\bvg\bva%
\llap{„}\alst{F}rýr⸗a maðr þér ęngi, Gunn·arr, \hld\ hęfir \alst{f}ull-vegit! &
Lítt sé⸗sk \alst{A}tli \hld\ \alst{ó}fu þïna; &
hann mun \alst{y}kkar \hld\ \alst{ǫ}nd síðarri &
ok \alst{ę́} vesa \hld\ \alst{a}fl it męira.\eva

\bvb “No one taunts thee, Guther; thou hast fully triumphed! \\
Little does Attle fear thy wrath; \\
he will breathe longer than you two, \\
and for ever be the greater power.\evb\evg


\bvg\bva%
\alst{S}ęgja mun’k þér, Gunn·arr, \hld\ ---\alst{s}jalfr vęitst gǫr⸗la--- &
hvé ér yðr \alst{s}nimma \hld\ til \alst{s}aka réðuð; &
varð’k⸗at \alst{e}k til \alst{u}ng \hld\ né \alst{o}f·þrungin &
\alst{f}ull-gǿdd \alst{f}é\emph{i} \hld\ ȧ \alst{f}lęti bróður.\eva

\bvb I will tell thee Guther---thou knowest it clearly thyself--- \\
how ye were too soon in your decisions. \\
I was not too young nor oppressed, \\
fully endowed with wealth upon my brother’s floor.\evb\evg


\bvg\bva%
Né ek \alst{v}ilda þat \hld\ at mik \alst{v}err ę́tti &
áðr ér \alst{G}júkungar \hld\ riðuð at \alst{g}arði &
\alst{þ}rír ȧ hęstum, \hld\ \edtrans{\alst{þ}jóð-konungar}{great kings}{\Bfootnote{An old poetic word; for etymology + distribution see \textlink{Gripisspa}[1]/2 n.}}, &
ęn \alst{þ}ęira fǫr \hld\ \alst{þ}ǫrf⸗gi vę́ri.\eva

\bvb Nor did I wish that a husband should have me \\
before ye Yivickings rode to the court, \\
three great kings upon your horses--- \\
but their coming was not a boon.\evb\evg


\bvg\bva%
\edtext{Þęim hét⸗umk þȧ}{\Bfootnote{There is certainly some corruption here seen by the unusual meter of this orphaned “line”, but the syntax and sense otherwise appears in order and no obvious emendation suggests itself.}} &
es með \alst{g}olli sat \hld\ ȧ \alst{G}rana bógum; &
vas⸗at hann ï \alst{au}gu \hld\ \alst{y}ðr of líkr, &
né ȧ \alst{ę}ngi hlut \hld\ at \alst{ȧ}·litum; &
\alst{þ}ó \alst{þ}ikki⸗sk ér \hld\ \alst{þ}jóð-konungar.\eva

\bvb Then I vowed myself to \emph{him} \\
who sat with the gold on the back of Grane. \\
In his eyes he was not like \emph{you} \\
nor in any part of his appearance, \\
though ye think yourselves great kings!\evb\evg


\bvg\bva%
Ok mér \alst{A}tli þat \hld\ \alst{ęi}nni sagði &
at \alst{h}vár⸗ki lét⸗sk \hld\ \alst{h}ǫfn of dęila, &
\alst{g}oll né jarðir, \hld\ nema \alst{g}efa⸗sk léta’k; &
ok \alst{ę}ngi hlut \hld\ \alst{au}ðins f\emph{éa}r, &
þȧ’s mér \alst{jó}ð-ungri \hld\ \alst{ęi}g\emph{u} sęldi &
ok mér \alst{jó}ð-ungri \hld\ \alst{a\emph{u}}ra talði.\eva

\bvb TODO. \\
TODO. \\
TODO. \\
TODO. \\
TODO. \\
TODO: Translation.\evb\evg


\bvg\bva%
Þȧ vas ȧ \alst{h}vǫrfun \hld\ \alst{h}ugr mïnn umb þat &
hvárt ek skylda \alst{v}ega \hld\ eða \alst{v}al fęlla, &
\alst{b}ǫll ï \alst{b}rynju, \hld\ umb \alst{b}róður sǫk. &
\alst{Þ}at myndi \alst{þ}ȧ \hld\ \alst{þ}jóð-kunnt vesa, &
\alst{m}ǫrgum \alst{m}anni \hld\ at \alst{m}unar stríði.\eva

\bvb TODO. \\
TODO. \\
TODO. \\
TODO. \\
TODO: Translation.\evb\evg


\bvg\bva%
Létum \alst{s}íga \hld\ \alst{s}átt-mǫ́l okkur; &
lék mér \alst{m}ęirr ï \alst{m}un \hld\ \alst{m}ęiðmar þiggja, &
\alst{b}auga rauða, \hld\ \alst{b}urar Sig·mundar, &
né ek \alst{a}nnars manns \hld\ \alst{au}ra vilda’k.\eva

\bvb TODO. \\
TODO. \\
TODO: Translation. \\
No other mans ounces of silver would I have.\evb\evg


\bvg\bva%
\alst{U}nna \alst{ęi}num \hld\ né \alst{ý}missum, &
bjó⸗at umb \alst{h}vęrfan \hld\ \alst{h}ug męn-Skǫgul; &
\alst{a}llt mun þat \alst{A}tli \hld\ \alst{ę}ptir finna &
es hann \alst{m}ïna spyrr \hld\ \alst{m}orð-fǫr gǫrva,\eva

\bvb TODO. \\
TODO. \\
TODO. \\
TODO: Translation.\evb\evg


\bvg\bva%
at \alst{þ}ęy⸗gi skal \hld\ \alst{þ}unn-gęð kona &
\alst{a}nnarrar ver \hld\ \alst{a}ldri lęiða; &
þȧ mun ȧ \alst{h}ęfndum \hld\ \alst{h}arma mïnna.“\eva

\bvb TODO. \\
TODO. \\
TODO: Translation.\evb\evg


\bvg\bva%
Upp ręis \alst{G}unn·arr, \hld\ \alst{g}ramr verðungar, &
ok umb \alst{h}als konu \hld\ \alst{h}ęndr of lagði; &
gingu \alst{a}llir, \hld\ ok þó \alst{ý}msir, &
af \alst{h}ęilum \alst{h}ug, \hld\ \alst{h}ana at lętja.\eva

\bvb TODO. \\
TODO. \\
TODO. \\
TODO: Translation.\evb\evg


\bvg\bva%
\alst{H}ratt af \alst{h}alsi \hld\ \alst{h}\emph{v}ęim þar sér; &
\alst{l}ét⸗a mann sik \alst{l}ętja \hld\ \alst{l}angrar gǫngu.\eva

\bvb TODO. \\
TODO: Translation.\evb\evg


\bvg\bva%
Nam hann sér \alst{H}ǫgna \hld\ \alst{h}vętja at ru̇num: &
\llap{„}\alst{S}ęggi vil’k alla \hld\ ï \alst{s}al ganga, &
\alst{þ}ïna með mïnum \hld\ ---nú ’s \alst{þ}ǫrf mikil--- &
vita ef \alst{m}ęini \hld\ \alst{m}orð-fǫr konu &
und’s af \alst{m}éli \hld\ ęnn \alst{m}ęin komi, &%TODO: check nasal vowel in méli
\alst{þ}ȧ lǫ́tum \alst{þ}ví \hld\ \alst{þ}arfar ráða.“\eva

\bvb TODO. \\
TODO. \\
TODO. \\
TODO. \\
TODO. \\
TODO: Translation.\evb\evg


\bvg\bva%
\alst{Ęi}nu því Hǫgni \hld\ \alst{a}nd·svǫr vęitti: &
\llap{„}\alst{L}ętj⸗a maðr hana \hld\ \alst{l}angrar gǫngu &
\edtext{þar’s hǫ̇n \alst{a}ptr·borin \hld\ \alst{a}ldri verði}{\lemma{þar’s hǫ̇n aptr·borin aldri verði ‘the long journey where she might never be born again’}\Bfootnote{The underlying implication is a bit unclear.  Is the suicide thought to prevent her from being reborn, or is the sense that she will be so happy with Siward in the afterlife that she her soul will not desire to return among the living?}}! &
Hǫ̇n \alst{k}rǫng of \alst{k}om⸗sk \hld\ fyr \alst{k}né móður, &
hǫ̇n \alst{ę́} borin \hld\ \alst{ȯ}·vilja til, &
\alst{m}ǫrgum \alst{m}anni \hld\ at \alst{m}óð-trega.“\eva

\bvb This one answer Hawn then granted: \\
“Let no man restrain her from the long journey \\
where she might never be born again! \\
Frail she arrived before her mother’s lap; \\
\emph{she} was always born to ill deeds, \\
to [cause] heart-grief for many a man.”\evb\evg


\bvg\bva%
Hvarf sér \alst{ȯ}·hróðugr \hld\ \alst{a}nd·spilli frȧ &
þar’s \alst{m}ǫrk \alst{m}ęnja \hld\ \alst{m}ęiðmum dęildi.\eva

\bvb The dishonoured man turned away from their conversation, \\
where the wood of neckrings \ken{woman} was dealing out treasures.\evb\evg


\bvg\bva%
\Ballnote{Byrnhild dons a golden byrnie and fatally wounds herself by disembowelement (see note to l. 3 \emph{miðlaði}).  There is a notable contradiction here, s Byrnhild cuts through the byrnie she has just donned.}%
Lęit hǫ̇n of \alst{a}lla \hld\ \alst{ęi}gu sïna, &
\alst{s}oltnar þýjar \hld\ ok \alst{s}al-konur; &
\alst{g}oll-brynju smó \hld\ ---vas⸗a \alst{g}ótt ï hug--- &
áðr sik \edtrans{\alst{m}iðlaði}{pierced her stomach}{\Bfootnote{\emph{miðla} usually means ‘to deal out, share’ but here has the unparallelled sense ‘cut something at the middle’.  The image is thus that Byrnhild plunges her sword into her belly (cf.~the Japanese \emph{hara-kiri} or ‘bowel-cutting’) and disembowels herself.}} \hld\ \alst{m}ę́kis ęggjum.\eva

\bvb She beheld throughout all her estate \\
the dead hand-maids and hall-women. \\
She donned a golden byrnie---there was no good in her heart--- \\
before she pierced her stomach with the short-sword’s edges.\evb\evg


\bvg\bva%
\alst{H}né við bólstri \hld\ \alst{h}ǫ̇n ȧ annan veg &
ok \alst{h}jǫr-unduð \hld\ \alst{h}ugði at rǫ́ðum:\eva

\bvb She knelt down against the bolster on her side \\
and, sword-wounded, thought of counsel:\evb\evg


\bvg\bva%
\llap{„}Nú skulu \alst{g}anga \hld\ þę́r’s \alst{g}oll vili &
ok \alst{m}inn\emph{i} því \hld\ at \alst{m}ér þiggja; &
ek gef \alst{h}vęrri \hld\ \edtrans{of \alst{h}roðit sigli}{a gold-enclosed necklace}{\Bfootnote{\emph{of hroðit} ‘(metal-)enclosed’ corresponding to OE \emph{ge·hroden}.  For \emph{sigli} ‘necklace’, only found thrice in the whole Norse corpus, see \textlink{Lokasenna}[20].}}, &
\alst{b}ók ok \alst{b}lę́ju, \hld\ \alst{b}jartar váðir.“\eva

\bvb “Now shall those who will have gold \\
and less things than it come and receive it from me. \\
I give every woman a gold-enclosed necklace, \\
a coverlet and bed-cover, [and] bright garments.”\evb\evg


\bvg\bva%
Þǫgðu all\emph{i}r, \hld\ hugðu at rǫ́ðum, &
ok \alst{a}ll\emph{i}r sęnn \hld\ \alst{a}nn·svǫr vęittu: &
\llap{„}\alst{Ǿ}rnar soltnar \hld\ munum \alst{ę}nn lifa, &
verða \alst{s}al-konur \hld\ \alst{s}ø̇mð at vinna.“\eva

\bvb All were silent; thought of counsel, \\
and all at once an answer granted: \\
“Enough women have died!  We will live on; \\
the women of the hall will do what is fitting.”\evb\evg


\bvg\bva%
\Ballnote{Byrnhild relents.  She will not force those who are reluctant to die with her, but they will receive greater honour in the afterlife if they follow her now than if they go on living.}%
Und’s af \alst{h}yggjandi \hld\ \alst{h}ǫr-skrýdd kona, &
\alst{u}ng at \alst{a}ldri, \hld\ \alst{o}rð viðr of kvað: &
\llap{„}Vil’k⸗at ek mann \alst{t}rauðan \hld\ né \alst{t}or·bø̇nan &
umb \alst{ȯ}ra sǫk \hld\ \alst{a}ldri tẏna,\eva

\bvb Until wisely the linen-clad woman \\
young of age quoth a word in reply: \\
“I wish not that ont reluctant or hardly asked \\
for our sake should lose his life,\evb\evg


\bvg\bva%
þó mun ȧ \alst{b}ęinum \hld\ \alst{b}rinna yðrum &
\alst{f}ę́ri ęyrir \hld\ þȧ’s ér \alst{f}ramm komið, &
nęitt \alst{M}ęnju góð, \hld\ \alst{m}ïn at vitja.\eva

\bvb yet will on your limbs burn fewer \\
ounces of silver when ye arrive \\
---nor any boons of Mene \ken{golden jewels}---to visit me.\evb\evg


\bvg\bva%
\Ballnote{Byrnhild speaks to Guther and makes a spae.}%
\alst{S}ęt⸗sk-tu niðr, Gunn·arr, \hld\ mun’k \alst{s}ęgja þér &
\alst{l}ífs ør·vę̇na \hld\ \alst{l}jósa brúði; &
mun⸗a \alst{y}ðvart far \hld\ \alst{a}llt ï sundi &
þó’tt \alst{e}k hafa \hld\ \alst{ǫ}ndu látit.\eva

\bvb Sit thee down Guther---I will tell thee \\
of a light bride whose hopes of life are o’er. \\
Your ship will not be all [sunk] in the sound, \\
though I have given up my breath of life.\evb\evg


\bvg\bva%
\alst{S}ǫ́tt munuð it Guð·ru̇n, \hld\ \alst{s}nemmr an hyggir; &
hęfir \alst{k}unn \alst{k}ona \hld\ við \alst{k}onung\emph{i} &
\edtrans{\alst{d}aprar minjar}{sorrowful remembrance}{\Bfootnote{Swanhild, by whom Guthrun will remember Siward.}} \hld\ at \alst{d}auðan ver.\eva

\bvb Thou and Guthrun will be reconciled sooner than thou thinkest; \\
the famous woman has a sorrowful remembrance \\
from the king after her husband’s death.\evb\evg


\bvg\bva%
Þar ’s \alst{m}ę́r borin, \hld\ \alst{m}óðir fǿðir; &
sú mun \alst{h}vítari \hld\ an inn \alst{h}ęiði dagr, &
\alst{S}van·hildr, vesa, \hld\ \edtrans{\alst{s}ólar gęisla}{sun-ray}{\Bfootnote{Swanhild is also likened to a \emph{sólar gęisli} in \textlink{Gudrunarhvot}[15].}}.\eva

\bvb There a maiden is born; her mother rears her. \\
She will be whiter than the clear day, \\
Swanhild, than a ray of the sun.\evb\evg


\bvg\bva%
\alst{G}efa munt \alst{G}uð·ru̇nu \hld\ \alst{g}óðra nǫkkurum, &
\alst{sk}ęyti \alst{sk}ǿða \hld\ \alst{sk}atna męngi; &
mun⸗at at \alst{v}ilja \hld\ \alst{v}er-sę́l gefin, &
hana mun \alst{A}tli \hld\ \alst{ęi}ga ganga, &
\edtrans{of \alst{b}orinn \alst{B}uðla \hld\ \alst{b}róðir mïnn}{the man born of Budle, my own brother}{\Bfootnote{Identically shared with \textlink{GudrunOne}[25]/3.}}.\eva

\bvb TODO. \\
TODO. \\
TODO. \\
Attle will go to take her to wife, \\
the man born of Budle, my own brother.\evb\evg


\bvg\bva%
\alst{M}args á’k \alst{m}inna⸗sk \hld\ hvé við \alst{m}ik fóru &
þȧ’s mik \alst{s}ára \hld\ \alst{s}vikna hǫfðuð; &
\alst{v}aðin at \alst{v}ilja \hld\ \alst{v}as’k meðan lifða’k.\eva

\bvb I have much to remember from how they treated me, \\
when they had sorely betrayed me. \\
I was deprived of willpower while I lived.\evb\evg


\bvg\bva%
Munt \alst{O}dd·ru̇nu \hld\ \alst{ęi}ga vilja &
ęn þik \alst{A}tli mun \hld\ \alst{ęi}gi láta; &
it munuð \alst{l}úta \hld\ ȧ \alst{l}aun saman, &
hǫ̇n mun þér \alst{u}nna \hld\ sęm \alst{e}k skylda’k, &
ef okkr \alst{g}óð of skǫp \hld\ \alst{g}ørði verða.\eva

\bvb Thou wilt want to take Guthrun to wife \\
but Attle will nowise let thee. \\
Ye two will embrace in secret. \\
She will love thee like \emph{I} should have \\
if good shapes had let it be so for us.\evb\evg


\bvg\bva%
Þik mun \alst{A}tli \hld\ \alst{i}llu bęita, &
munt ï \alst{ø}ngan \hld\ \alst{o}rm-garð lagiðr.\eva

\bvb Thee will Attle afflict with evil; \\
thou wilt be laid in the narrow snake-pit.\evb\evg


\bvg\bva%
\alst{Þ}at mun ok verða \hld\ \alst{þ}ví⸗gi⸗t lęngra &
at \alst{A}tli mun \hld\ \alst{ǫ}ndu tẏna, &
\alst{s}ę́lu \alst{s}ïnni, \hld\ ok \alst{s}ofa lífi. &
Því’t hǫ̇num \alst{G}uð·ru̇n \hld\ \alst{g}rýmir ȧ bęð &
\alst{s}nǫrpum ęggjum \hld\ af \alst{s}ǫ́rum hug.\eva

\bvb TODO. \\
TODO. \\
TODO. \\
TODO. \\
TODO: Translation.\evb\evg


\bvg\bva%
\alst{S}ø̇mri vę́ri Guð·ru̇n, \hld\ \alst{s}ystir okkur &
frum-ver sïnum \hld\ [...] &
ef hęn\emph{n}i \alst{g}ę́fi \hld\ \alst{g}óðra ráð &
eða \alst{ę́}tti hǫ̇n hug \hld\ \alst{o}ss of líkan.\eva

\bvb It would be more fitting for Guthrun, our sister, \\
{[to follow]} her first husband {[in death]}, \\
if she had the counsel of good men, \\
or she had a heart like ours.\evb\evg


\bvg\bva%
\alst{Ȯ}·\alst{ǫ}rt mę́li’k nú \hld\ ęn hǫ̇n \alst{ęi}gi mun &
umb \alst{ȯ}ra sǫk \hld\ \alst{a}ldri tẏna; &
\edtrans{\alst{h}ana munu \alst{h}ęfja \hld\ \alst{h}ǫ́var bǫ́rur}{She will be lifted by billows high}{\Bfootnote{Clearly related to \textlink{Gudrunarhvot}[13]/3, with which the b-verse is identical.}} &
til \alst{Jȯ}n·akrs \hld\ \alst{ó}ðal-torfu.\eva

\bvb Now I speak unhastily, but she will not \\
lose her life for our sake. \\
She will be lifted by billows high \\
to Enacker’s ethel-turf.\evb\evg


\bvg\bva%
{[...]} eru ï vǫr-úðum \hld\ Jȯn·akrs sonum; &
mun hǫ̇n \alst{S}van·hildi \hld\ \alst{s}ęnda af landi, &
\alst{s}ïna męy \hld\ ok \alst{S}ig·\emph{va}rðar.\eva

\bvb TODO. \\
she will send Swanhild away from the land, \\
her daughter and Siward’s.\evb\evg


\bvg\bva%
\Ballnote{Byrnhild predicts Swanhild’s death.  Cf.~\textlink{Gudrunarhvot}[P1].}%
Hana munu \alst{b}íta \hld\ \alst{B}ikka rǫ́ð &
því’t \alst{Jǫ}rmun·rekkr \hld\ \alst{ȯ}·þarft lifir; &
þȧ’s \alst{ǫ}ll farin \hld\ \alst{ę́}tt Sig·\emph{va}rðar, &
eru \alst{G}uð·ru̇nar \hld\ \alst{g}rǿti at flęiri.\eva

\bvb Her will Beck’s counsels bite, \\
for Ermenric lives wickedly. \\
Then is all the lineage of Siward destroyed; \\
for Guthrun there are yet more causes for weeping.\evb\evg


\bvg\bva%
\alst{B}iðja mun’k þik \hld\ \alst{b}ø̇nar einnar, &
sú mun ï \alst{h}ęimi \hld\ \alst{h}inst bø̇n vesa: &
Lát-tu svá \alst{b}ręiða \hld\ \alst{b}org ȧ vęlli &
at undir \alst{o}ss \alst{ǫ}llum \hld\ \alst{ja}fn-rúmt sé\emph{i}, &
þęim es \alst{s}ultu \hld\ með \alst{S}ig·\emph{vę}rði.\eva

\bvb I will ask thee a single ask; \\
in the Home it will be my very last ask. \\
Let the stronghold upon the plain be so broad \\
that under us all there be even room--- \\
us who perished with Siward.\evb\evg


\bvg\bva%
\alst{T}jaldi þar umb þȧ borg \hld\ \alst{t}jǫldum ok skjǫldum, &
\alst{v}ala-ript \alst{v}el fǫ́ð \hld\ ok \edtrans{\alst{v}ala}{horses}{\Bfootnote{A rare word.  TODO.}} męngi; &
bręnni mér inn \alst{h}u̇nska \hld\ ȧ \alst{h}lið aðra.\eva

\bvb Let it be canopied there around the stronghold with canopies and shields, \\
Welsh canvas well painted and a multitude of horses. \\
Let me be burned by the Hunnish man’s side.\evb\evg


\bvg\bva%
Bręnni inum \alst{h}u̇nska \hld\ ȧ \alst{h}lið aðra &
\alst{m}ïna þjȯna, \hld\ \alst{m}ęnjum gǫfga, &
\alst{t}vȧ at hǫfðum \hld\ ok \alst{t}vęir haukar; &
þȧ’s \alst{ǫ}llu skipt \hld\ til \alst{ja}fnaðar.\eva

\bvb Let be burned by the Hunnish man’s other side \\
my servants noble of neckrings, \\
two by his head, and two hawks; \\
then all is arranged evenly.\evb\evg


\bvg\bva%
\Ballnote{Cf.~\textlink{Brot}[20].}%
Liggi okkar ęnn ï \alst{m}illi \hld\ \edtrans{\alst{m}almr hring-variðr}{ring-adorned metal \ken{sword}}{\Bfootnote{A ring-sword, long obsolete by the time of \Sigurdskamma.  See \textlink{HelgakvidaHjorvardssonar}[9]/1 n.}}, &
\alst{ę}gg-hvasst \emph{\alst{é}a}rn, \hld\ svá \alst{ę}ndr lagit &
þȧ’s vit \alst{b}ę́ði \hld\ \alst{b}ęð ęinn stigum &
ok \alst{h}étum þȧ \hld\ \alst{h}jȯna nafni.\eva

\bvb Let once more lie between us the ring-adorned metal \ken{sword}, \\
the edge-sharp iron, once again laid \\
like when we both mounted a single bed \\
and then were called by a wedded couple’s name.\evb\evg


\bvg\bva%
\alst{H}rynja \alst{h}ǫ̇num þȧ \hld\ ȧ \alst{h}ę̇l þęy⸗gi &
\edtrans{\alst{h}lunn-blik \alst{h}allar}{roller-gleams of the hall}{\Bfootnote{TODO, proverbial?}}, \hld\ \alst{h}ring\emph{a} litkuð &
ef hǫ̇num \alst{f}ylgir \hld\ \alst{f}ęrð mïn heðan; &
þęy⸗gi mun \alst{v}ǫ̇r fǫr \hld\ \alst{au}m⸗lig vesa,\eva

\bvb Then will nowise the roller-gleams \ken{doors?} of the hall \\
crash down on the heels of the colourer of ring-swords \ken{warrior}, \\
if he is followed by my departure hence. \\
Nowise will our journey be wretched,\evb\evg


\bvg\bva%
því’t hǫ̇num \alst{f}ylgja \hld\ \alst{f}imm ambóttir, &
\alst{á}tta þjȯnar, \hld\ \alst{ø}ðlum góðir, &
\alst{f}ostr-man mitt \hld\ ok \alst{f}aðerni &
þat’s \alst{B}uðli gaf \hld\ \alst{b}arni sïnu.\eva

\bvb for he is followed by five handmaids, \\
eight servants of good pedigree, \\
my foster-daughter, and the paternity \\
which Budle gave unto his child [me].\evb\evg


\bvg\bva%
\alst{M}argt sagða ek, \hld\ \alst{m}ynda’k flęira &
es mér \alst{m}ęir \alst{m}jǫtuðr \hld\ \alst{m}ál-rúm gę́fi; &
\alst{ó}mun þverr, \hld\ \alst{u}ndir svella, &
\alst{s}att ęitt \alst{s}agða’k--- \hld\ \alst{s}vá mun’k láta!“\eva

\bvb Many things I said, I would [say] yet more \\
if the Measurer gave me greater room to speak. \\
My voice wanes, my wounds swell, \\
truth alone I have spoken---so will I die!”\evb\evg

\sectionline
